\section{Positive mixed Hankel determinants and finite arithmetic reconstruction}
\label{sec:positive-mixed-hankel}

Throughout this section put $K=\Eexp$.

Put $H_0=0$, $H_m=\sum_{j=1}^m1/j$, and define
\begin{align}
 G_m
 &=\int_0^\infty x^me^{-x}(x-\log x)\,dx
   =m!(\gamma+m+1-H_m),
 \label{eq:mixed-G-moment}\\
 I_{k+1}
 &=\int_0^1t^k\frac{e^{1-1/t}}t\,dt
   =\int_0^\infty\frac{e^{-x}}{(1+x)^{k+1}}\,dx.
 \label{eq:mixed-I-moment}
\end{align}
Integration by parts gives $I_{m+1}=(1-I_m)/m$.  Hence
\begin{equation}\label{eq:mixed-I-arithmetic}
 k!I_{k+1}=A_k+(-1)^k\delta,
 \qquad A_0=0,
 \qquad A_k=(k-1)!-A_{k-1}\in\mathbb Z.
\end{equation}
For $d\geq1$ set
\begin{equation}\label{eq:mixed-Hankel-definition}
 M_k^{(d)}=G_{k+d-1}I_{k+1},
 \qquad
 \mathcal P_d(\gamma,\delta)
 =\det\bigl(M_{i+j}^{(d)}\bigr)_{0\leq i,j<d}.
\end{equation}

\begin{theorem}[Positive full-degree mixed determinant]
\label{thm:positive-full-degree-mixed-determinant}
For every $d\geq1$,
\[
 \mathcal P_d(\gamma,\delta)\in\mathbb Q[\gamma,\delta],
 \qquad \mathcal P_d(\gamma,\delta)>0.
\]
Its total degree is exactly $2d$, and its top mixed coefficient is
\begin{equation}\label{eq:mixed-Hankel-top-coefficient}
 [\gamma^d\delta^d]\mathcal P_d
 =(-1)^{d(d-1)/2}\bigl((d-1)!\bigr)^d\ne0.
\end{equation}
More precisely, it has the Heine representation
\begin{align}
 \mathcal P_d
 &=\frac1{d!}
 \int_{((0,\infty)\times(0,1))^d}
 \prod_{1\leq i<j\leq d}(x_it_i-x_jt_j)^2
 \notag\\[-1mm]
 &\qquad\times
 \prod_{\ell=1}^d
 x_\ell^{d-1}e^{-x_\ell}(x_\ell-\log x_\ell)
 \frac{e^{1-1/t_\ell}}{t_\ell}
 \,dx_\ell\,dt_\ell.
 \label{eq:mixed-Hankel-Heine}
\end{align}
\end{theorem}

\begin{proof}
The product variable $y=xt$ under the positive measure
\[
 x^{d-1}e^{-x}(x-\log x)
 \frac{e^{1-1/t}}t\,dx\,dt
\]
on $(0,\infty)\times(0,1)$ has kth moment $M_k^{(d)}$.
Andreief's identity gives \eqref{eq:mixed-Hankel-Heine}.  Since
$x-\log x\geq1$ and the product measure is non-atomic with infinite
support, the Vandermonde is nonzero almost everywhere off a null set;
the integral is therefore strictly positive.

Equations \eqref{eq:mixed-G-moment} and
\eqref{eq:mixed-I-arithmetic} prove polynomial membership and the upper
bound $2d$ for the total degree.  The coefficient of $\gamma\delta$ in
$M_k^{(d)}$ is
\[
 (-1)^k\frac{(k+d-1)!}{k!}
 =(-1)^k(k+1)_{d-1},
\]
where $(x)_n$ is the rising factorial.  Consequently the coefficient of
$\gamma^d\delta^d$ in the determinant equals
\[
 \det\bigl((-1)^{i+j}(i+j+1)_{d-1}\bigr)_{0\leq i,j<d}.
\]
Factoring the row and column signs has no net effect.  For
$p(k)=(k+1)_{d-1}$, apply the unitriangular forward-difference transform
to both rows and columns.  The resulting matrix has entries
$\mathsf D^{i+j}p(0)$; it vanishes past the anti-diagonal, while every
anti-diagonal entry is
$\mathsf D^{d-1}p(0)=(d-1)!$.  Its determinant is therefore the right
side of \eqref{eq:mixed-Hankel-top-coefficient}.  This coefficient is
nonzero, so the total degree is exactly $2d$.
\end{proof}

\begin{proposition}[Finite reconstruction from the first two mixed determinants]
\label{prop:mixed-hankel-reconstruction}
Put $U=\mathcal P_1$ and $V=\mathcal P_2$ for the determinants in
Theorem~\ref{thm:positive-full-degree-mixed-determinant}.  Then the extension
\[
 K(\gamma,\delta)/K(U,V)
\]
is algebraic of degree at most four.  More explicitly,
\[
 U=(\gamma+1)\delta,
\]
and $\delta$ is a root of
\begin{equation}\label{eq:mixed-hankel-quartic}
 2T^4+(U-4)T^3
 +(2U^2-16U+2V+2)T^2
 +(-16U^2+8U)T+8U^2=0,
\end{equation}
after which $\gamma=U/\delta-1$.
Consequently
\[
 \operatorname{trdeg}_K K(U,V)
 =\operatorname{trdeg}_K K(\gamma,\delta).
\]
In particular, at least one of $\mathcal P_1,\mathcal P_2$ is
transcendental over $K$, and the explicit number
\[
 \mathcal P_1+i\mathcal P_2
\]
is transcendental over $K$.  Moreover $\mathcal P_1,\mathcal P_2$ are
algebraically independent over $K$ if and only if $\gamma,\delta$ are.
\end{proposition}

\begin{proof}
The formulas for $G_m$ and $I_m$ give
$\mathcal P_1=(\gamma+1)\delta$.  The explicit second determinant is
\[
 2\mathcal P_2=-18+36\delta-5\delta^2-24\gamma
 +48\gamma\delta-5\gamma\delta^2-8\gamma^2
 +16\gamma^2\delta-2\gamma^2\delta^2.
\]
Substituting $\gamma=U/\delta-1$ and clearing $\delta^2$ gives
\eqref{eq:mixed-hankel-quartic}.  Since $\delta>0$, the displayed
recovery of $\gamma$ is valid.  This proves finiteness and equality of
transcendence degrees over $K$.

The pure exponential-base theorem makes
$\eta=\Ein(1)=\gamma+\delta/e$ transcendental over $K$.  Since $e\in K$,
the two numbers $\gamma,\delta$ cannot both be algebraic over $K$.
Therefore $U,V$ cannot both be algebraic over $K$.  The field $K$ is
stable under complex conjugation.  If $U+iV$ were algebraic over $K$,
then its conjugate would be algebraic over $K$, and hence so would its
real and imaginary parts $U,V$, a contradiction.  The final equivalence
is the case of transcendence degree two.
\end{proof}
